%!TeX encoding = UTF-8
%!TEX program = xelatex
% ============================================================
% 《波谱分析》课程作业论文
% 题目：红外光谱与拉曼光谱在纳米材料表征中的比较与联用决策
% 版式沿用 X-Ray/作业二/课程论文/main.tex（色带题头 + 校徽、楷体、蓝色标题层级）
% 编译：XeLaTeX + Biber，运行本目录 complie.bat main
% ============================================================
\documentclass[zihao=5,a4paper,twoside]{ctexart}

% ==================== 页面 ====================
\usepackage[top=2.4cm,bottom=2.5cm,left=1.5cm,right=1.4cm,footskip=0.55cm]{geometry}

% ==================== 字体 ====================
\setCJKmainfont{KaiTi}[AutoFakeBold=2.0]
\setmainfont{Times New Roman}
\usepackage{indentfirst}
\setlength{\parindent}{2em}
\setlength{\parskip}{0pt}
\raggedbottom

% ==================== 颜色 ====================
\usepackage{xcolor}
\definecolor{secblue}{HTML}{365F91}
\definecolor{subblue}{HTML}{243F60}
\definecolor{bandblue}{HTML}{9CB8C7}

% ==================== 图表与数学 ====================
\usepackage{graphicx}
\usepackage{float}
\usepackage{booktabs}
\usepackage{array}
\usepackage{tabularx}
\usepackage{amsmath,amssymb}
\usepackage{bm}
\usepackage{caption}
\usepackage[labelfont=bf]{subcaption}
\usepackage{enumitem}
\usepackage{tikz}
\usetikzlibrary{arrows.meta,positioning,calc}

% ==================== 参考文献 ====================
\usepackage[backend=biber,style=gb7714-2015]{biblatex}
\addbibresource{ref.bib}
\setcounter{biburlnumpenalty}{10000}
\setcounter{biburlucpenalty}{10000}
\setcounter{biburllcpenalty}{10000}

\usepackage[hidelinks]{hyperref}
\hypersetup{
  pdftitle={红外光谱与拉曼光谱在纳米材料表征中的比较与联用决策},
  pdfauthor={刘佳鹏},
  pdfsubject={《波谱分析》课程论文},
  pdfkeywords={红外光谱, 拉曼光谱, TiO2, 羧酸根配位, 选择定则}
}

% ==================== 编号（按章编号，如 图 2-1）====================
\numberwithin{figure}{section}
\numberwithin{table}{section}
\numberwithin{equation}{section}
\renewcommand{\thefigure}{\thesection-\arabic{figure}}
\renewcommand{\thetable}{\thesection-\arabic{table}}
\renewcommand{\theequation}{\thesection-\arabic{equation}}

\captionsetup{font=small,labelfont=bf,labelsep=quad}

% ==================== 标题样式 ====================
\ctexset{
  section = {
    format     = {\color{secblue}\bfseries\fontsize{13pt}{16pt}\selectfont},
    number     = {\thesection.},
    aftername  = {\hspace{0.4em}},
    beforeskip = 15pt plus 3pt minus 2pt,
    afterskip  = 7pt plus 2pt,
  },
  subsection = {
    format     = {\color{secblue}\bfseries\fontsize{12pt}{15pt}\selectfont},
    aftername  = {\hspace{0.4em}},
    beforeskip = 12pt plus 2pt,
    afterskip  = 5pt plus 1pt,
  },
  subsubsection = {
    format     = {\color{subblue}\bfseries\fontsize{11pt}{14pt}\selectfont},
    aftername  = {\hspace{0.4em}},
    beforeskip = 10pt plus 2pt,
    afterskip  = 4pt plus 1pt,
  },
  paragraph = {
    format     = {\color{subblue}\bfseries\fontsize{11pt}{14pt}\selectfont},
    runin      = false,
    beforeskip = 10pt plus 2pt,
    afterskip  = 4pt plus 1pt,
  },
}

% ==================== 页脚 ====================
\usepackage{fancyhdr}
\pagestyle{fancy}
\fancyhf{}
\renewcommand{\headrulewidth}{0pt}
\fancyfoot[RO]{\fontsize{8pt}{10pt}\selectfont |\enspace\thepage}
\fancyfoot[LE]{\fontsize{8pt}{10pt}\selectfont \thepage\enspace|}
\fancypagestyle{plain}{%
  \fancyhf{}%
  \renewcommand{\headrulewidth}{0pt}%
  \fancyfoot[RO]{\fontsize{8pt}{10pt}\selectfont |\enspace\thepage}%
}

\begin{document}
\linespread{1.2}\selectfont

% ==================== 首页题头 ====================
\hspace*{0.25cm}%
\begin{minipage}[c]{15.75cm}
  \setlength{\fboxsep}{0pt}%
  \colorbox{bandblue}{%
    \parbox[c][0.73cm][c]{15.75cm}{%
      \hspace{0.06cm}\textcolor{white}{\fontsize{16pt}{16pt}\selectfont\songti 《波谱分析》课程论文}}}
\end{minipage}\hfill
\begin{minipage}[c]{1.4cm}
  \centering\includegraphics[height=1.27cm]{assets/MUC.jpg}
\end{minipage}

\vspace{0.75cm}

% ==================== 论文题目 ====================
{\centering\bfseries\fontsize{14pt}{20pt}\selectfont 红外光谱与拉曼光谱在纳米材料表征中的比较与联用决策\par}
\vspace{0.25cm}
{\centering\songti\fontsize{11pt}{15pt}\selectfont ——以油酸包覆 TiO$_2$ 纳米颗粒的表面配位与晶相分析为例\par}

\vspace{0.55cm}

\hspace*{3.3cm}{\songti\fontsize{12pt}{16pt}\selectfont 刘佳鹏，24012639，纳米二班}

\vspace{0.45cm}

% ==================== 摘要 ====================
% ==================== 摘要 ====================
{\bfseries 摘要}\quad
红外与拉曼都探测分子振动，但物理来源与选择定则不同。本文以油酸包覆 TiO$_2$ 纳米颗粒为例，围绕“油酸以何种羧酸根配位、其与晶相有何关系”比较两种技术。本人认为应以红外判配位，拉曼补晶相与分布，联用可建立“配位—晶相—位置”的关联，并讨论水、荧光、制样与激光损伤等限制。

\vspace{4pt}

\par\noindent{\heiti 关键词：}红外光谱；拉曼光谱；TiO$_2$ 纳米颗粒；羧酸根配位；选择定则；联用


% ==================== 正文 ====================
% ==================== 1 材料体系与表征问题 ====================
\section{材料体系与表征问题}

TiO$_2$ 纳米颗粒广泛用于光催化与防晒，其晶相、粒径和包覆层决定分散稳定性与活性\cite{sacco2023}。油酸是常用包覆剂，羧酸根可能以质子化羧酸吸附，也可能以单齿、双齿螯合或桥联方式与 Ti 位点配位\cite{lv2013}，配位方式影响包覆层致密性、分散性及相变。本文的问题因此定为：油酸在 TiO$_2$ 表面以何种羧酸根配位？其与晶相组成及相稳定性有何关系？

% ==================== 2 物理来源、选择定则与互斥规则 ====================
\section{物理来源、选择定则与互斥规则}

\subsection{物理来源}

两者都源于分子振动能级跃迁，但物理过程不同。红外属吸收过程：振动必须改变分子偶极矩（$\Delta\mu\neq0$），光子能量正好等于振动能级差时被吸收。拉曼属散射过程：激光把分子激发到虚能级后几乎立即返回，若返回能级与出发时不同，散射光便产生波数差（拉曼位移）；强度取决于极化率变化（$\Delta\alpha\neq0$），常温下斯托克斯线远强于反斯托克斯线（图~\ref{fig:energy}）。

\begin{figure}[H]
  \centering
  \begin{tikzpicture}[font=\footnotesize,>=Latex,line width=0.6pt]
    % (a) 红外吸收
    \node at (1.2,4.15) {(a) 红外吸收};
    \draw (0,0) -- (2.4,0);
    \draw (0,1.6) -- (2.4,1.6);
    \draw[dashed] (0,3.2) -- (2.4,3.2);
    \node[left] at (0,0) {$v=0$};
    \node[left] at (0,1.6) {$v=1$};
    \node[left] at (0,3.2) {虚能级};
    \draw[->,thick] (1.2,0) -- (1.2,1.6);
    \node[right] at (1.25,0.8) {$h\nu_0=\Delta E$};
    % (b) 拉曼散射
    \node at (7.5,4.15) {(b) 拉曼散射};
    \draw (4.5,0) -- (10.5,0);
    \draw (4.5,1.6) -- (10.5,1.6);
    \draw[dashed] (4.5,3.2) -- (10.5,3.2);
    \node[left] at (4.5,0) {$v=0$};
    \node[left] at (4.5,1.6) {$v=1$};
    \node[left] at (4.5,3.2) {虚能级};
    % 瑞利
    \draw[->] (5.4,0) -- (5.4,3.2);
    \draw[->] (5.6,3.2) -- (5.6,0);
    \node[below] at (5.5,-0.05) {瑞利};
    % 斯托克斯
    \draw[->] (7.0,0) -- (7.0,3.2);
    \draw[->] (7.2,3.2) -- (7.2,1.6);
    \node[below] at (7.1,-0.05) {斯托克斯};
    \node[right] at (7.25,2.4) {$h(\nu_0-\nu_{\mathrm{vib}})$};
    % 反斯托克斯
    \draw[->] (8.9,1.6) -- (8.9,3.2);
    \draw[->] (9.1,3.2) -- (9.1,0);
    \node[below] at (9.0,-0.05) {反斯托克斯};
    \node[right] at (9.15,2.4) {$h(\nu_0+\nu_{\mathrm{vib}})$};
    \node[right] at (5.65,2.2) {$h\nu_0$};
  \end{tikzpicture}
  \caption{红外吸收与拉曼散射的能级过程。}
  \label{fig:energy}
\end{figure}

\subsection{选择定则与互斥规则的适用条件}
\noindent\textbf{1. 选择定则。}群论中，IR 活性要求振动模式的不可约表示与偶极矩分量（$x$、$y$、$z$）的不可约表示之一相同；Raman 活性则要求该不可约表示与极化率张量分量（$x^2$、$y^2$、$z^2$、$xy$、$yz$、$zx$）的不可约表示之一相同。

\noindent\textbf{2. 互斥规则。}它只在有对称中心时成立：中心对称群的振动只能属于 g（对反演对称）或 u（反对称）两类，偶极矩属 u 型、极化率属 g 型，同一振动不能同时具有两种活性。

\noindent\textbf{3. 适用条件。}结构确有对称中心，如 CO$_2$（$D_{\infty h}$）、锐钛矿（$I4_1/amd$）与金红石（$P4_2/mnm$）；比较对象为同一振动模式；结构为理想周期，表面、缺陷与吸附会降低局域对称性，使禁戒模弱激活\cite{taudul2023}。

CO$_2$ 对称伸缩 1333 cm$^{-1}$ 只见于拉曼、反对称伸缩 2349 cm$^{-1}$ 只见于红外\cite{larkin2017}；TiO$_2$ 的 E$_g$ 144/447 cm$^{-1}$ 属 g 模，u 模落在 400 cm$^{-1}$ 以下远红外。油酸无对称中心，两谱均可观察，但极性 COO$^-$ 的偶极矩变化大、极化率变化小，红外强度远高于拉曼\cite{ning2014}。所以红外主要判断配位，拉曼主要判断晶格与晶相。

% ==================== 3 技术选择：优先红外还是拉曼 ====================
\section{技术选择}

\subsection{优先红外光谱}

在技术选择上，红外具有明显优势，原因如下：1. 判据直接。COO$^-$ 的 $\nu_{\mathrm{as}}$ 与 $\nu_{\mathrm{s}}$ 在红外里都是强峰，拉曼截面却小得多。2. 判据可量化。配位方式由 $\Delta\nu=\nu_{\mathrm{as}}-\nu_{\mathrm{s}}$ 判定：明显小于 110 cm$^{-1}$ 为双齿螯合，接近或大于 200 cm$^{-1}$ 为单齿，桥联介于 140--200 cm$^{-1}$\cite{deacon1980}。3. 干扰可识别。1710 cm$^{-1}$ 的 C=O 若仍然很强，表明样品里还有游离油酸\cite{piao2026ir}。与油酸、油酸钠及模型化合物做差谱与分峰，三种配位可区分（表~\ref{tab:compare}）。

\begin{table}[H]
  \centering
  \caption{红外与拉曼的比较}
  \label{tab:compare}
  \small
  \begin{tabularx}{\textwidth}{@{}lXX@{}}
    \toprule
    项目 & 红外（ATR-FTIR） & 拉曼（785 nm） \\
    \midrule
    物理过程 & 吸收，$\Delta\mu\neq0$ & 非弹性散射，$\Delta\alpha\neq0$ \\
    本体系关键信息 & COO$^-$ 的 $\nu_{\mathrm{as}}$/$\nu_{\mathrm{s}}$、$\Delta\nu$、C=O、O—H & 144/447 cm$^{-1}$、C—H/C=C \\
    水与干扰 & 水干扰强；KBr 吸潮 & 水峰弱；荧光强、信号弱 \\
    空间分辨率 & 1.5--10 $\mu$m & 250--500 nm \\
    \bottomrule
  \end{tabularx}
\end{table}

\subsection{拉曼光谱能够补充的信息}

拉曼光谱能够补充以下信息：1. 晶相组成。锐钛矿 144 cm$^{-1}$ 与金红石 447 cm$^{-1}$ 峰形尖锐，峰强比或 PLS 模型都能定量，5\%--95\% 内偏差可优于 2\%\cite{sacco2023}，红外低波数吸收难以做到。2. 空间分布。油酸的 C=C 与 C—H 在拉曼里也有响应，样品可在水相直接测量并做 mapping\cite{kurouski2020}。不足是 COO$^-$ 信号弱，$\Delta\nu$ 的判据不如红外可靠。

% ==================== 4 单技术方案与联用决策 ====================
\section{单技术方案与联用决策}

\subsection{只能选一种时的方案}

若只保留一台仪器，方案仍应以红外为骨架，因为本问题问的是配位方式，红外对羧酸根最敏感。方案不能照搬前文，需改以下几处：KBr 压片换成 ATR 或 DRIFTS，避开吸潮与散射；1520--1560 cm$^{-1}$ 的重叠峰用差谱与分峰拆开；配位判据用纯相 TiO$_2$、油酸钠与油酸标样重新标定。低波数 Ti—O—Ti 只能给晶相半定量\cite{piao2026ir}。若只剩拉曼，则用 785 nm、低功率、长积分，先测相组成，再由 C—H/C=C 判断包覆层；而配位方式只能作倾向性判断，不能给出确切的结论。

\subsection{联用能多知道什么}

两种谱叠在同一视场上，才能表明清楚配位与晶相是否相关。红外只表明油酸怎样结合，拉曼只表明哪里是哪种相；把二者对齐后，若某一晶相的 $\Delta\nu$ 或构象不同，就可判断油酸是否优先稳定该晶相。这种“配位—晶相—位置”的对应，单一技术不能表明清楚\cite{kurouski2020}。

% ==================== 5 实际测量限制与应对 ====================
\section{实际测量限制与应对}

\noindent\textbf{水的干扰：}红外最棘手，O—H 伸缩（3400 cm$^{-1}$）与 1640 cm$^{-1}$ 弯曲会掩盖 COO$^-$ 信息；应对是样品干燥成膜、ATR 保持干燥并扣背景\cite{piao2026ir}。

\noindent\textbf{样品制备：}KBr 压片易吸潮，还会引入散射与基线倾斜；改用 ATR 或 DRIFTS 可少制样，但压力过大会压坏包覆层。

\noindent\textbf{荧光背景：}油酸杂质或缺陷会产生荧光，可能淹没 144 cm$^{-1}$ 以外的弱峰；785 nm 比 532 nm 荧光低，但信号按 $\lambda^{-4}$ 下降，可用长波长与基线拟合抑制。

\noindent\textbf{信号强度：}拉曼截面极小，COO$^-$ 尤弱，需高数值孔径物镜与多次累加；SERS 虽能增强信号，但基底会改变表面化学\cite{wang2021}。

\noindent\textbf{激光损伤：}功率过高会使油酸脱附、局部升温，甚至诱发锐钛矿→金红石相变或 Ti$^{4+}$→Ti$^{3+}$ 还原\cite{taudul2023}；表面功率宜控制在 1 mW 以下并散焦。

\noindent\textbf{空间分辨率：}常规红外受衍射限制约 1.5--10 $\mu$m，拉曼约 250--500 nm；分辨单颗粒需 AFM-IR/TERS（10--50 nm 与 3--5 nm），但设备复杂且可能扰动包覆层\cite{kurouski2020}；定量还需纯相标样或 PLS 校准\cite{sacco2023}。

% ==================== 6 结论 ====================
\section{结　论}

回到最初的问题：红外表明“油酸怎样结合”，拉曼表明“什么相、多大、如何分布”。只选一种时，红外能表明清楚配位方式，却不能表明清楚配位与晶相的关系；两者叠合，才能把配位、晶相与位置串起。


% ==================== 参考文献 ====================
\printbibliography[title=参考文献]

\end{document}
